\documentclass[11pt]{amsart}
\usepackage[a4paper,margin=28mm]{geometry}
\usepackage{amsmath,amssymb}
\usepackage[colorlinks=true,linkcolor=blue,citecolor=blue,urlcolor=blue]{hyperref}
\hypersetup{pdftitle={Total integrability of a Minkowski billiard in its unit ball},
pdfauthor={Kaiwen Jin}}
\newtheorem{theorem}{Theorem}[section]
\newtheorem{lemma}[theorem]{Lemma}
\newtheorem{proposition}[theorem]{Proposition}
\theoremstyle{remark}
\newtheorem{remark}[theorem]{Remark}
\newcommand{\R}{\mathbb R}
\newcommand{\area}{\operatorname{area}}
\newcommand{\tr}{\operatorname{tr}}
\newcommand{\dd}{\,d}
\title{Total integrability of a Minkowski billiard in its unit ball}
\author{Kaiwen Jin}
\email{yui.ui.siki@gmail.com}
\date{October 15, 2026}
\subjclass[2020]{37E40, 52A10, 53A15}
\keywords{Minkowski billiards, total integrability, convex curves}
\begin{document}
\begin{abstract}
Let a smooth strongly convex norm on the plane define both a billiard
table and the length of its chords. We prove that if the entire billiard
phase annulus is foliated by continuous essential invariant curves, then
the norm comes from an inner product. No finite symmetry beyond central
symmetry is assumed. The proof retains the support point of each chord
as an independent variable, converts length variations into alternating
area variations, and averages over the two resulting phase annuli.
A centro-affine boundary parameter and the Blaschke--Santal\'o inequality
then force the unit ball to be an ellipse.
\end{abstract}
\maketitle

\section{Introduction}\label{sec:intro}
In a Minkowski billiard, straight chords are measured by a norm and the
reflection law is obtained by stationarity of their total norm length.
When the table is the unit ball of that same norm, one may ask whether
a foliation of all reflected trajectories detects an inner-product norm.
This is the rigidity question attributed to
Y.~Ostrover in \cite[Section 9.5, Question 1]{BM}.

Let $N:\R^2\to[0,\infty)$ be a norm, smooth away from zero, with
$D^2(N^2)$ positive definite there, and put $K=\{N\leq1\}$.
Let $\gamma(s)$ parametrize $\partial K$ by positive $N$-arclength,
with period $T$. The billiard map acts on ordered distinct boundary
points, represented by $s<t<s+T$, and has generating function
\[
L(s,t)=N(\gamma(t)-\gamma(s)).
\]
Its consecutive parameters satisfy
\[
L_2(s_{j-1},s_j)+L_1(s_j,s_{j+1})=0.
\]
We use \emph{total integrability} to mean that this entire phase annulus
is foliated by continuous invariant simple closed curves winding once
around it.

Bialy and Tsodikovich proved rigidity under an additional finite linear
symmetry of order at least three
\cite[Theorem 5]{BT}; their formulation uses a $C^1$ foliation.
Their argument uses periodic invariant curves to produce progressively
more symmetries. A general norm ball has central symmetry, of order two,
which does not supply that initial higher-order symmetry.
The following theorem proves the rigidity implication for continuous
invariant foliations without the additional assumption.

\begin{theorem}\label{thm:main}
For every norm $N$ as above, if its billiard in $K=\{N\leq1\}$ is totally
integrable, then $K$ is an origin-centered ellipse. Equivalently, there
exists a real symmetric positive-definite matrix $Q$ such that
\[
N(v)=\sqrt{v^TQv}\qquad\text{for every }v\in\R^2.
\]
\end{theorem}

The proof uses an area variational description in position and momentum.
This description and its relation to symplectic billiards were studied
by Albers, Chavez Caliz, and Tabachnikov \cite[Sections 2--3]{ACT}.
For a single symplectic billiard table, total-integrability rigidity
was proved by Baracco and Bernardi \cite[Theorem 1.1]{BB}.
Here the position curve and the rotated polar momentum curve need not
coincide, so a single-table rigidity theorem does not directly apply.
Instead, each norm chord is the maximum of an area expression over its
support point. This maximum transfers finite variational inequalities
to an alternating sequence on the two curves. Averaging requires both
alternating phase annuli.

The remaining estimate is geometric. We choose a boundary parameter for
which the signed area integral of the derivative curve equals that of
the original curve. Its period is at most $2\pi$ by the
Blaschke--Santal\'o inequality. The averaged variational inequality and
an energy lower bound then force equality in that inequality.
The assumption concerns the whole phase annulus; integrability only
near its boundary is outside the theorem.
Section~\ref{sec:area} develops the alternating variational inequality,
and Section~\ref{sec:affine} proves the geometric estimates and the theorem.

\section{From norm length to alternating area}\label{sec:area}

Throughout, set
\[
\omega(x,y)=x_1y_2-x_2y_1,\qquad J(x_1,x_2)=(-x_2,x_1).
\]
For a convex body containing the origin, its polar is
$C^*=\{p:p\cdot x\leq1\text{ for every }x\in C\}$.
Let
\begin{equation}\label{eq:bodies}
P=-JK^*,\qquad A=\area(K),\qquad B=\area(P)=\area(K^*).
\end{equation}
Both boundaries are smooth, have positive curvature, and are centrally
symmetric. All boundary parameters increase counterclockwise.

\subsection{Support points and finite variations}

Keeping the support point of a chord as a separate variable turns its norm
length into a sum of two signed areas.
\begin{lemma}\label{lem:support}
For distinct $x,z\in\partial K$,
\begin{equation}\label{eq:support}
N(z-x)=\max_{y\in\partial P}
\{\omega(x,y)+\omega(y,z)\}.
\end{equation}
The maximizer is unique and has strictly negative second derivative in
any regular boundary parameter of $P$.
\end{lemma}
\begin{proof}
The expression being maximized is $(z-x)\cdot Jy$, and $JP=K^*$.
The support function of $K^*$ is $N$. For instance, convexity and
homogeneity show that $DN(z-x)\in K^*$ and
$DN(z-x)\cdot(z-x)=N(z-x)$.
Positive curvature of the polar gives uniqueness and nondegeneracy:
at the supporting point a nonzero outward supporting functional has
strictly negative second derivative along the boundary.
\end{proof}

We first record the variational consequence of total integrability.
For a positive regular parameter $\gamma$ on $\partial K$, put
$w=\gamma(t)-\gamma(s)$. Convexity and homogeneity give
$D^2N(w)w=0$. Moreover $D^2N(w)$ is positive on every direction
transverse to $w$. Otherwise, subtracting a multiple of $w$ would
give a nonzero $u$ with $DN(w)u=0$ and
$D^2(N^2)(w)[u,u]=0$. Thus, in the plane,
\[
D^2N(w)[u,v]=\lambda(w)\omega(w,u)\omega(w,v),
\qquad \lambda(w)>0.
\]
At the two endpoints of a chord,
$\omega(\gamma'(s),w)>0$ and $\omega(\gamma'(t),w)<0$.
Consequently the norm-length generating function has
\begin{equation}\label{eq:twist}
L_{12}(s,t)=-D^2N(w)[\gamma'(s),\gamma'(t)]>0.
\end{equation}
In the original $N$-arclength parameter, $-L_1(s,t)$ decreases from
$1$ to $-1$ as $t$ runs from $s$ to $s+T$. These coordinates
identify the ordered-point annulus with the usual symplectic twist
cylinder.

For this sign convention, every orbit on a continuous essential invariant
curve is locally maximizing for the action with fixed endpoints; see
\cite[Section 3, paragraph after Theorem 3.1, p.~5]{Bialy}.
The hypothesis of Theorem~\ref{thm:main} therefore makes the Hessian
of every finite length action negative semidefinite. This property
is unchanged by a smooth increasing boundary reparametrization:
at a stationary action its Hessian changes by diagonal congruence.

For even vertices $q_i\in\partial K$ and odd vertices $p_i\in\partial P$,
consider
\begin{equation}\label{eq:alternating}
\sum_i\{\omega(q_i,p_i)+\omega(p_i,q_{i+1})\}.
\end{equation}
At the supporting $p_i$, maximizing each odd variable in
\eqref{eq:alternating} gives exactly the original length action.
The odd-variable Hessian is diagonal negative definite by
Lemma~\ref{lem:support}. If the full finite Hessian is
$\begin{pmatrix}U&V\\V^T&W\end{pmatrix}$, its Schur complement
$U-VW^{-1}V^T$ is the Hessian of that length action. Completing the square,
\[
\begin{pmatrix}\xi\\\eta\end{pmatrix}^{\!T}
\begin{pmatrix}U&V\\V^T&W\end{pmatrix}
\begin{pmatrix}\xi\\\eta\end{pmatrix}
=\xi^T(U-VW^{-1}V^T)\xi
+(\eta+W^{-1}V^T\xi)^TW(\eta+W^{-1}V^T\xi)\leq0.
\]
Every finitely supported alternating variation can be enclosed between
fixed even endpoints. Hence all its second variations are nonpositive.

\subsection{The two half maps}

Let $c(s)$ and $d(t)$ be positive regular parametrizations of
$\partial K$ and $\partial P$. Define the two annuli
\[
\mathcal D_{KP}=\{(s,t):\omega(c'(s),d'(t))>0\},\qquad
\mathcal D_{PK}=\{(t,s):\omega(d'(t),c'(s))>0\}.
\]
For $(q,p)=(c(s),d(t))\in\mathcal D_{KP}$, let $z$ be the other
intersection of $\partial K$ with the ray $q+u d'(t)$, $u>0$, and set
$S(q,p)=(p,z)$. The determinant condition means the ray enters $K$.
It exits transversely at $z$, where
$\omega(d'(t),c'(z))>0$; here $c'(z)$ denotes the positive tangent at $z$.
Thus the image belongs to $\mathcal D_{PK}$.
On this component define $S(p,z)=(z,p_+)$ using the other intersection
of $P$ with the ray from $p$ in the positive tangent direction at $z$.
Reverse chords give the inverse half maps. Transversality makes them smooth.
We regard $S$ as one invertible map of the disjoint union
$\mathcal D=\mathcal D_{KP}\sqcup\mathcal D_{PK}$.

The point $p$ is the maximizer for $N(z-q)$ in \eqref{eq:support}.
Indeed the outward normal to $P$ at $p$ is a positive multiple of
$-Jd'(t)$, and $z-q$ is a positive multiple of $d'(t)$.
Conversely every ordered distinct $q,z$ determines exactly such a $p$.
At the next even vertex, the stationarity equation
\[
\omega(c'(z),p_+-p)=0
\]
is, by differentiation of \eqref{eq:support}, the norm-length reflection
equation. Its outgoing branch is the other intersection $p_+\ne p$:
the equal intersection would continue the incoming chord out of $K$.
Therefore $S^2$ on $\mathcal D_{KP}$ is conjugate to the original
Minkowski billiard. Every edge in either component lies on an alternating
lift of an original orbit. The finite second-variation conclusion above
holds on the whole union.

On each component use $H(u,v)=\omega(u,v)$, differentiating in its
respective boundary parameters, and write $b=H_{12}>0$.
The measure
\begin{equation}\label{eq:measure}
\dd\mu=b\dd s\dd t
\end{equation}
on each component is finite and the combined measure on $\mathcal D$
is invariant under $S$. To see this directly, differentiate
$H_2(s,t)+H_1(t,s_+)=0$ and wedge with $\dd t$:
\[
H_{12}(s,t)\dd s\wedge\dd t
=H_{12}(t,s_+)\dd t\wedge\dd s_+.
\]
Both densities are positive on their corresponding annuli.
Finiteness follows from bounded derivatives on the compact parameter tori.

\subsection{An averaged inequality}

At $n$ consecutive interior vertices choose variation $1$, with variation
$0$ at the two external endpoints. The nonpositive second variation is
\[
\sum_{j=1}^n\{H_{22}(u_{j-1},u_j)+H_{11}(u_j,u_{j+1})\}
+2\sum_{j=1}^{n-1}H_{12}(u_j,u_{j+1})\leq0.
\]
Integrate over the initial edge in the invariant union $\mathcal D$.
There are $n$ diagonal terms and $n-1$ adjacent pairs, so invariance gives
\[
n\int_{\mathcal D}(H_{11}+H_{22})\dd\mu
+2(n-1)\int_{\mathcal D}H_{12}\dd\mu\leq0.
\]
All integrands are bounded and the measure is finite. Dividing by $n$
and letting $n\to\infty$ yields
\begin{equation}\label{eq:average}
\int_{\mathcal D}(H_{11}+2H_{12}+H_{22})H_{12}\dd s\dd t\leq0.
\end{equation}

Choose the parameters to respect antipodes:
$c(s+\ell_c/2)=-c(s)$ and $d(t+\ell_d/2)=-d(t)$.
Such parameters will be constructed in Lemma~\ref{lem:parameter}.
Replacing one endpoint by its antipode reverses every second derivative
of $H$. It leaves the quadratic integrand in \eqref{eq:average}
unchanged and exchanges the positive- and negative-determinant domains.
The determinant-zero integrand vanishes. Thus each component integral is
half of its full product-torus integral. Interchanging the two bodies
negates $H$, interchanges $H_{11},H_{22}$, and leaves that product
unchanged. The two full integrals agree. Consequently
\eqref{eq:average} is precisely
\begin{equation}\label{eq:fulltorus}
\iint\{\omega(c'',d)+2\omega(c',d')+\omega(c,d'')\}
\omega(c',d')\dd s\dd t\leq0,
\end{equation}
with both parameters running through full periods.

We express this in terms of the energy matrices
\begin{equation}\label{eq:energies}
\begin{split}
M_c&=\int c'c'^T\dd s,\qquad M_d=\int d'd'^T\dd t,\\
C_c&=\int\omega(c',c'')\dd s,\qquad
C_d=\int\omega(d',d'')\dd t,\\
D&=\iint\omega(c',d')^2\dd s\dd t
=\tr(M_cJM_dJ^T).
\end{split}
\end{equation}
Recall that $A=\area(K)$ and $B=\area(P)=\area(K^*)$.
Periodic integration by parts and the area formula give
$\int d d'^T\dd t=-BJ$ and $\int c c'^T\dd s=-AJ$.
At fixed $s$ the first term in \eqref{eq:fulltorus} therefore integrates
to $-B\omega(c',c'')$; interchanging the bodies gives the other term.
It follows that
\begin{equation}\label{eq:hopf}
2D\leq BC_c+AC_d.
\end{equation}

\section{Affine estimates and the equality case}\label{sec:affine}

To use \eqref{eq:hopf}, we choose parameters for which each derivative
integral $C_c,C_d$ is twice the area of its original body.
The same parameters have a universal period bound.

\begin{lemma}\label{lem:parameter}
Let $C=-C$ be a smooth planar convex body with positive curvature
and area $a>0$. There is a positive boundary parameter $\sigma$,
of period $\Omega_C\leq2\pi$, such that
\begin{equation}\label{eq:derivativearea}
\int_0^{\Omega_C}\omega(q_\sigma,q_{\sigma\sigma})\dd\sigma=2a.
\end{equation}
Antipodal points differ by $\Omega_C/2$ in this parameter.
\end{lemma}
\begin{proof}
First use the centro-affine area parameter $r$, defined by
$\omega(q,q_r)=1$. It exists because the origin is strictly interior;
its full period is $2a$. Differentiation and positive curvature give
\[
q_{rr}=-k(r)q,\qquad k(r)>0.
\]
The polar boundary is $v=-Jq_r$. Indeed $v$ is the outward supporting
covector at $q$ and $v\cdot q=1$. Since $v_r=kJq$,
$\omega(v,v_r)=k$, so, with $a^*=\area(C^*)$,
\[
\int k(r)\dd r=2a^*.
\]
Set
\[
\dd\sigma=\sqrt{k(r)}\dd r,\qquad
\Omega_C=\int\sqrt{k(r)}\dd r.
\]
Cauchy--Schwarz and the symmetric Blaschke--Santalo inequality
\cite[Theorem 1]{BK} give
\begin{equation}\label{eq:period}
\Omega_C^2\leq4aa^*\leq4\pi^2.
\end{equation}
Direct differentiation gives
\[
q_\sigma=\frac{q_r}{\sqrt{k}},\qquad
q_{\sigma\sigma}=-q-\frac{k_r}{2k^2}q_r,\qquad
\omega(q_\sigma,q_{\sigma\sigma})=\frac1{\sqrt{k}}.
\]
Integrating proves \eqref{eq:derivativearea}.
Central symmetry gives $q(r+a)=-q(r)$ and $k(r+a)=k(r)$,
so antipodes are half a $\sigma$ period apart.
\end{proof}

\begin{lemma}\label{lem:energy}
Let $q$ be a smooth regular simple closed planar curve, enclosing area
$a>0$, in a parameter of period $\ell$. Then its energy matrix
$M=\int q'q'^T$ satisfies
\begin{equation}\label{eq:detbound}
\sqrt{\det M}\geq\frac{2\pi a}{\ell}.
\end{equation}
\end{lemma}
\begin{proof}
The matrix $M$ is positive definite: a vector in its kernel would make
one component of $q$ constant and put the curve in a line.
The area-preserving linear map
$T=(\det M)^{1/4}M^{-1/2}$ satisfies
$TMT^T=\sqrt{\det M}\,I$. If $l$ is the Euclidean perimeter of $Tq$,
Cauchy--Schwarz and ordinary isoperimetry
\cite[Section 1, equation (1.5)]{Osserman} imply
\[
4\pi a\leq l^2
\leq\ell\int|Tq'|^2
=2\ell\sqrt{\det M}.
\]
\end{proof}

\begin{proof}[Proof of Theorem~\ref{thm:main}]
Parametrize both boundaries by Lemma~\ref{lem:parameter}.
In \eqref{eq:energies} this gives $C_c=2A$, $C_d=2B$,
with periods $\Omega_K,\Omega_P\leq2\pi$.
The variational inequality \eqref{eq:hopf} becomes
\begin{equation}\label{eq:upper}
D\leq2AB.
\end{equation}
On the other hand, for positive-definite $2$ by $2$ matrices $U,V$,
\[
\tr(UV)\geq2\sqrt{\det U\det V},
\]
by the arithmetic--geometric mean inequality for the eigenvalues
of $U^{1/2}VU^{1/2}$.
Lemma~\ref{lem:energy} applied separately to $c,d$ therefore gives
\begin{equation}\label{eq:lower}
D=\tr(M_cJM_dJ^T)
\geq2\sqrt{\det M_c\det M_d}
\geq\frac{8\pi^2AB}{\Omega_K\Omega_P}
\geq2AB.
\end{equation}
The separate linear maps used to prove \eqref{eq:detbound} establish
scalar determinant bounds; the mixed trace in \eqref{eq:lower}
is computed in the original coordinates.

Since $A,B>0$, \eqref{eq:upper}--\eqref{eq:lower} force
$\Omega_K\Omega_P=4\pi^2$ and hence $\Omega_K=\Omega_P=2\pi$.
Since $B=\area(K^*)$, applying \eqref{eq:period} to $C=K$ yields
$4\pi^2\leq4AB\leq4\pi^2$.
Equality in the symmetric Blaschke--Santalo inequality makes $K$
an origin-centered ellipse. Its gauge is
$N(v)=\sqrt{v^TQv}$ for a real symmetric positive-definite matrix $Q$,
including $v=0$.
\end{proof}

\subsection*{AI use}
OpenAI Codex was used to develop candidate arguments, retrieve and check
references, and draft and revise the manuscript.
Independent mathematical review is requested.

\begin{thebibliography}{9}
\bibitem{ACT}
P. Albers, A. Chavez Caliz, and S. Tabachnikov,
\emph{Symplectic billiards as Minkowski billiards},
arXiv:2607.05986v1 (2026).
\url{https://arxiv.org/abs/2607.05986v1}.

\bibitem{BB}
L. Baracco and O. Bernardi,
\emph{Totally integrable symplectic billiards are ellipses},
arXiv:2305.19701v1 (2023).
\url{https://arxiv.org/abs/2305.19701v1}.

\bibitem{Bialy}
M. Bialy,
\emph{Hopf rigidity for convex billiards on the Hemisphere and Hyperbolic Plane},
arXiv:1205.3873v2 (2012).
\url{https://arxiv.org/abs/1205.3873v2}.

\bibitem{BM}
M. Bialy and A. E. Mironov,
\emph{Integrable Billiards and Related Topics},
arXiv:2510.03790v2 (2025).
\url{https://arxiv.org/abs/2510.03790v2}.

\bibitem{BT}
M. Bialy and D. Tsodikovich,
\emph{Billiard tables with rotational symmetry},
arXiv:2106.06956v2 (2021).
\url{https://arxiv.org/abs/2106.06956v2}.

\bibitem{BK}
G. Bianchi and M. Kelly,
\emph{The Blaschke--Santal\'o Inequality},
arXiv:1312.0244v3 (2014).
\url{https://arxiv.org/abs/1312.0244v3}.

\bibitem{Osserman}
R. Osserman,
\emph{The isoperimetric inequality},
Bull. Amer. Math. Soc. \textbf{84} (1978), no.~6, 1182--1238.
\url{https://doi.org/10.1090/S0002-9904-1978-14553-4}.
\end{thebibliography}
\end{document}
